\vspace*{-3ex}
\subsection{Impacts and Relations of the Proposed Research}
\label{impact}
\vspace*{-1ex}

{\bf Our research is aimed at advancing image and video database
retrieval and indexing by addressing a board range of critical problems 
in such enterprise, including the design of features, adaptive 
distance measures, and learning algorithms.} 

\paragraph{Invariant features}
The proposed framework of invariants ameliorates many difficulties
encountered in computing invariant image features. Because the
basis functions used (wavelets \cite{Daubechies88,Mallat89a},
short-time Fourier analysis \cite{Gabor46,Rabiner78}, and splines
\cite{Farin83,Rogers90}) have compact support, the proposed technique
tolerates a certain degree of self-occlusion unlike the invariants based on
global analysis \cite{Hu61,Teh88}.  Furthermore, it does not require
estimating any derivatives such as in the invariants using local analysis
\cite{Rivlin95,Weiss93b}.  Whence, the new method is more robust.  

The marriage of global models and local invariants allow us to
address the recognition and retrieval of objects of complex structures
and articulating parts. The use of local invariants allows component
objects to be recognized in a way insensitive to changes in lighting
and viewpoint.  The use of global models is beneficial in that
variation of an object's appearance due to articulated motion can be
accommodated.  By augmenting local invariants with global models, we
achieve efficient recognition of complex patterns insensitive to many
environmental changes.

Finally, {\bf the solutions to the {\em invariant, intra-class retrieval}
problem can have significant practical
implications} and will facilitate the design of retrieval subsystems
for real-world image databases.

\paragraph{Perceptual Similarity}
\label{sec-related-dpf}

Extensive work has been done on mapping color spectrums and textures to
the perceptual space \cite{B99,MWNS01,TMY78,W95,Jiawei98}.  Until
recently, little has been done on finding a good perceptual distance
function for aggregating individual perceptual features (work that has
been done will be discussed shortly).

Most researchers quantify the difference between two feature vectors by
their Minkowski distance or its variants.  However, in
\cite{Stricker95}, it was shown that for image retrieval the Minkowski
distance results in many false negatives and thus yields low recall.
Rubner et al. \cite{Rtg98} propose Earth Mover's Distance (EMD) for
image retrieval and claim marginal improvement (about $10$ percentile
in recall) over the Minkowski distances. However, their results may not
be conclusive, since only a few queries were tested, and those were on
small datasets.  The weighted Minkowski function introduces feature
weights into the distance function \cite{PMO99,RH00}. The shortcoming
of this metric is that images similar to the query image cannot be
captured by just one static set of weights.  For instance, both a
down-sampled and a cropped image are perceptually similar to the
original image.  But the similarity of the down-sampled image is
primarily captured in the color features (texture features are changed
when an image is down-sampled), whereas the similarity of the cropped
image may be captured in the texture features (color histograms are
changed when an image is cropped).  Using a static feature subset to
find both down-sampled and cropped images will not work effectively
because it cannot adapt to such changes.

Some other distance functions such as the quadratic-form distances
proposed by \cite{img11, ISF98}, have been shown to perform marginally
better than the Minkowski distance function.  Recently, the fractional
distance function \cite{Agg-ICDT01} has been shown to outperform both
L1 and L2 distance functions.  Our work shows that even better
functions can be derived, and the function that we derive (presented in
Section~\ref{sec-dpf}) not only works well but can also be interpreted
by perceptual principles in cognitive science.

Feature reduction has been shown to be a useful method for improving
the quality of similarity searches in the text domain.  The data is
transformed to a lower-dimensional space where most information is
preserved, by using methods such as Principle Component Analysis
(PCA).  Although PCA-like feature reduction has been widely used, the
work of C. Aggarwal \cite{Agg-PODS01} points out that the effectiveness
of feature reduction can be domain-dependent and data-dependent.
Moreover, for image retrieval, we have not found PCA very useful for
two reasons.  First, the perceptual features employed to characterize
images (such as color, texture, and shape) are sufficiently independent
of one another.  This agrees with C. Aggarwal's observation that when
the dimensions are relatively uncorrelated to one another, dimension
reduction can be counter-productive.  Second, and more importantly,
image perception is subjective, and hence feature reduction in the
image domain must be conducted in a user-dependent way, not in a
data-dependent way.  Our approach is first to use learning methods such
as $\MEGA$ and $\SVMActive$ to perform implicit feature transformation
in a user-dependent and query-dependent manner.  After a more effective
feature space is found, our dynamic partial distance function (DPF)
performs dynamic feature reduction for each pair of images. The
similarity of different pairs of images is measured in different
subspaces.  We explain in Section~\ref{sec-dpf} why this method works
effectively both empirically and theoretically.

\paragraph{Active Learning Algorithms}

Machine learning and relevance feedback techniques 
have been proposed to learn and to refine query concepts.
The problem is that most traditional techniques
require a large number of training instances \cite{Bish98,KV94,Mit97,ZH01}, 
and they require seeding a query with ``good'' 
examples \cite{JW97,PCM99,WFSP00}.
Unfortunately, in many practical scenarios, 
a learning algorithm must work with
a scarcity of training data
and a limited amount of training time.

Of late, ensemble techniques such as {\em bagging} \cite{Bre96},
{\em arcing} \cite{Bre98},
and {\em boosting} \cite{Grove+Schuurmans:AAAI98,Schapire+al:ICML97} 
have been proposed to improve learning
accuracy for decision trees and neural networks.
These ensemble schemes enjoy success in improving
classification accuracy
through bias or variance reduction,
but they do not help reduce the number of samples and time required
to learn a query concept.  In fact, most ensemble schemes
actually increase learning time because they
introduce learning redundancy in order to improve prediction
accuracy \cite{DB-95,Grove+Schuurmans:AAAI98,JH97,MM-98}.

To reduce the number of required samples, researchers
in the machine learning community have conducted 
several studies of active learning \cite{Cohn+al:1996,CFH97}
for classification.  Active learning can be modeled formally 
as follows: Given a dataset $S$ consisting of 
an unlabeled subset $U$ and a labeled subset $X$, an active
learner $L$ has two components: $f$ and $s$. The $f$ component
is a classifier that is trained on the current set of
labeled data $X$. The second component $s$ is the sampling
function that, given a current labeled set $X$, decides which
subset $u$ in $U$ to select to query the user.
The active learner returns a new $f$ after each round
of relevance feedback.  The sampling techniques
employed by the active learner determine the selection of 
the next batch of unlabeled instances to be labeled by the user.

The {\em query by committee} (QBC)
algorithm~\cite{Freund+al:ML97,Seung+al:COLT92} is
a representative active learning scheme.  QBC uses a distribution
over all possible classifiers and attempts greedily to reduce the
entropy of this distribution. This general purpose algorithm has been
applied in a number of domains (although, to our knowledge,
not to the image retrieval domain) using classifiers (such as Naive
Bayes classifiers~\cite{Dagan+Engelson:ICML95,McCallum+Nigam:ICML98})
for which specifying and sampling classifiers from a distribution is
natural.  Probabilistic models such as the Naive Bayes classifier
provide interpretable results and principled ways to incorporate prior
knowledge. However, they typically do not
perform as well as discriminative methods such as
SVMs~\cite{Joachims:ECML97,Dumais+al:CIKM98}, especially when
the amount of training data is scarce.
(See our SVM-based approach in Section~\ref{sec-svm}.)
For image retrieval where a query concept is typical non-linear\footnote{
A query such as ``animals'', ``women'', and 
``european architecture'' does not reside contiguously in the
space formed by the image features.}, our $\MEGA$ 
and $\SVMActive$ with kernel mapping
provide more flexible and accurate concept modeling.

%Two computational challenges in applying active
%learning in general and the QBC algorithm in
%particular to a learning task are 
%\begin{itemize} 
%\item Finding useful committee members quickly, and
%\item Finding informative unlabeled instances efficiently.
%\end{itemize} 

Specifically for image retrieval, the PicHunter 
system \cite{PicHunter96a,PicHunter96b,PicHunter97,CMM00} 
uses Bayesian
prediction to infer the goal image, based upon users' input.
Mathematically, the goal of PicHunter is 
to find a single goal point in the feature
space (e.g., a particular flower image), whereas our goal is
to hunt down all points that match a query concept
(e.g., the entire flower category,
which consists of flowers of different colors, shapes, and
textures, and against different backgrounds).  Note that
the points matching a target concept can be scattered all over 
the feature space. To find these points
quickly with few hints, 
our learning algorithms must deal with many daunting challenges.

The study of \cite{PMO99} puts relevance feedback techniques
proposed by the Information Retrieval (IR) 
into three categories: {\it query reweighting},
{\it query point movement} and {\it query expansion}.
\begin{jlist}
\item {\it Query reweighting} and {\it query point movement}
\cite{ISF98,ORCWMH98,ORCMH97,PCM99,RHOM98}.
Both query reweighting and query point movement
use nearest-neighbor sampling: They return top ranked
objects to be examined by the user and then refine the query
based on the user's feedback. If the initial query example
is good and the query concept is convex in the feature space
\cite{ISF98,WFSP00}, this nearest-neighbor sampling approach works
fine. Unfortunately, most users do not have a good example to
start a query, and most image-query concepts are non-convex.
%Refining a search around bad examples
%is analogous to trying to
%find oranges in the middle of an apple orchard by refining one's
%search to a few rows of apple trees at a time.
%It will take a long time to find oranges (the desired result).

\item {\it Query expansion} \cite{PMO99,WFSP00}.
The {\it query expansion} approach can be
regarded as a multiple-instances sampling approach.
The samples of a subsequent round are selected from the
neighborhood (not necessarily the nearest ones) of the
positive-labeled instances of the previous round. The study
of \cite{PMO99} shows that query expansion
achieves only a slim margin of improvement (about $10\%$
in precision/recall) over query point movement.
\end{jlist}

All traditional relevance feedback methods
require seeding the methods with ``good''
positive examples \cite{img11,img13,HVO99,MM96,SC,img12,img9},
and most methods do not use negative-labeled instances effectively.
For instance, sunset images must be supplied 
as examples in order to search for sunset pictures.
As we stated in Section~\ref{sec-overview},
finding good examples is the job of a search engine itself.
Our learning methods effectively use negative-labeled
instances to corner the positive instances, and 
select the most informative unlabeled instances to query
user feedback.  Because of the effective use of negative
and unlabeled instances, our method can learn a query
concept much faster and more accurately than the traditional
relevance-feedback methods.

%In short, the relevance feedback schemes that
%require seeding a query, and ineffective methods of
%sample selection, are not very useful at this stage.  Our proposed
%$\MEGA$ algorithm does not require seeding a query.
%(See Section~\ref{sec-mega}.)  Both $\MEGA$ and $\SVMActive$ employ 
%intelligent sampling methods that maximize
%the usefulness of each sample, and thereby
%achieve much better learning performance than earlier methods.



